%==============================================================================
% Section 5 · Method: EntroDiff (~2 pages)
%
% 三个组件 + 损失家族（与 path_A_method_skeleton §3 完全对齐）：
%   §5.1 Viscosity-matched noise schedule
%   §5.2 BV-aware score parameterization (核心 novelty)
%   §5.3 Loss family
%   §5.4 Reverse-time sampler with Godunov-form guidance
%==============================================================================
\section{Method: EntroDiff}
\label{sec:method}

% ----- §5.1 Viscosity-matched schedule -----
\subsection{Viscosity-matched noise schedule}
\label{sec:method:schedule}

\todo{Motivation: from \eqref{eq:score-burgers} the effective viscosity of the score-level Burgers is $g(\tau)^{2}/2$. Aligning it with the physical viscosity of the target PDE makes the two Burgers share the same viscous profile. We choose}

\begin{equation}
\label{eq:visc-matched}
\sigma^{2}(\tau) = 2\,\physvis\,\tau, \qquad \tau \in [0, T_{d}].
\end{equation}

% ----- §5.2 BV-aware parameterization -----
\subsection{BV-aware score parameterization}
\label{sec:method:bv-aware}

\todo{Key novelty. Score Shocks Proposition 5.4 gives the exact tanh form of $\score$ in the interfacial layer. We embed this form directly into the network architecture, decomposing the score into a smooth background gradient plus a tanh-shaped jump:}

\begin{equation}
\label{eq:bv-aware}
\sth(u, \tau) = \nabla \phi^{\mathrm{sm}}_{\theta}(u, \tau)
              + \frac{\kappa_{\theta}(u, \tau)}{2} \tanh\!\left( \frac{\phi^{\mathrm{sh}}_{\theta}(u, \tau)}{2} \right) \nabla \phi^{\mathrm{sh}}_{\theta}(u, \tau),
\end{equation}

\todo{where $\phi^{\mathrm{sm}}$ is a smooth background potential, $\phi^{\mathrm{sh}}$ is a signed-distance-to-shock function (zero on the shock set), and $\kappa$ is the local jump magnitude (constrained by Rankine–Hugoniot). This parameterization severs the $\exp(\amp)$ amplification source identified in Score Shocks Theorem 6.3.}

% ----- §5.3 Loss family -----
\subsection{Loss family}
\label{sec:method:loss}

\todo{Define the four loss components (Table~\ref{tab:loss-family}) and the final composite loss:}

\begin{equation}
\label{eq:loss-final}
\mathcal{L} = \Ldsm + \lambda_{\mathrm{ent}}\,\Rent(\Dth) + \lambda_{\mathrm{BV}}\,\TV(\Dth)
            + \lambda_{\mathrm{Burg}}\,\bigl\| \dtau \sth + 2\,\sth\,\dx \sth - \partial_{uu} \sth \bigr\|^{2}.
\end{equation}

\begin{table}[t]
    \caption{\todo{Loss-family ablation table. Key take-away: each regularizer addresses a specific theoretical role — entropy regularizer enforces Kruzhkov admissibility, BV regularizer keeps $\hat u$ in the entropy-solution function class, Burgers-consistency regularizer pushes $\sth$ onto the score Burgers manifold.}}
    \label{tab:loss-family}
    \centering
    \begin{tabular}{llll}
        \toprule
        Loss & Form & Theoretical role & Empirical role \\
        \midrule
        $\Ldsm$ & DSM & baseline regression target & generalization \\
        $\Lent$ & Kruzhkov entropy regularizer & physical-solution selector & shock-stability \\
        $\Lbv$ & total-variation penalty & keeps $\hat u \in \BV$ & oscillation control \\
        $\Lburg$ & score-Burgers consistency & manifold of admissible scores & robustness near shock \\
        \bottomrule
    \end{tabular}
\end{table}

% ----- §5.4 Sampler -----
\subsection{Reverse-time sampler with Godunov-form guidance}
\label{sec:method:sampler}

\todo{Use EDM-style probability-flow ODE plus DPS-style guidance, but replace the central-difference PDE residual (used by DiffusionPDE) with a Godunov-form entropy-consistent flux difference. This single change is what makes the guidance term well-defined across shocks.}

\begin{equation}
\label{eq:reverse-ode}
\frac{du}{d\tau} = -\sigtau\,\dot\sigma(\tau)\,\sth(u, \tau)
                  - \zeta_{\mathrm{obs}}\,\nabla_{u}\,\mathcal{L}_{\mathrm{obs}}(u)
                  - \zeta_{\mathrm{PDE}}\,\nabla_{u}\,\Lpde^{\mathrm{Godunov}}(u).
\end{equation}
